\chapter{A Timeline Is Not History Itself}
\section{Time on a Copper Coin}
First, pick up something small: a copper coin.

Not a valuable antique, just an ordinary Kangxi tongbao, a round coin with a square hole cast during the Kangxi reign of the Qing dynasty. So many survive that a little over ten yuan can buy one at an antiques stall. Hold it between your fingers. The four characters on its face read from top to bottom and right to left: Kang, xi, tong, bao.

Now a seemingly unnecessary question: in what year was this coin made?

You will find you cannot answer, and not because you studied history poorly. Kangxi is not a year but an emperor's reign title. Kangxi year one was 1662 CE; year sixty-one was 1722. Sixty years separate them, and the coin might have been cast in any of those years. It actually carries a complete system of reckoning time: time marked by an emperor's title rather than a simple numerical sequence. Saying ``This is Kangxi year thirty'' was as natural to a Qing person as saying ``This is 2025'' is to you.

Now consider where our readily spoken ``2025 CE'' comes from. Its count begins with the year of Jesus's birth, although most scholars today think even that starting point is off by several years. A monk devised this dating system only in the sixth century CE, and it took more than a thousand years to become the world's default.

In other words, ``What numbered year is this?'' has never been a natural fact. It is an invention. Nature gives us day and night, the waxing and waning moon, and the seasons. Humans decide to count years from a starting point. Who counts, where they begin, and how they divide the sequence shape the entire river of history in our minds.

That is this chapter's subject. The timeline you memorized in school is not history itself. A timeline is a map, and a map is always smaller than the territory. We will examine how that map is drawn, what it leaves out, and how it might be drawn better.

\section{The Night a City Fell}
Now come with me into a scene.

Time: the evening of May 28, 1453 CE. Place: Constantinople, capital of the Byzantine Empire, a city more than eleven hundred years old. On today's maps, it is Istanbul.

Defenders stand on the walls, perhaps only about seven thousand, including sailors and mercenaries who have come from Venice and Genoa. On the plain outside are tens of thousands of soldiers under the Ottoman sultan Mehmed II, their campfires stretching to the horizon. Outside, too, stand dozens of cannon among the world's largest at the time. They have battered these walls all spring. The defenders patch breaches by day and keep patching by night. They have done so for nearly two months.

The people know this is the last night. Scouts report that the sultan has ordered a general assault before dawn. That evening, the last emperor, Constantine XI, makes a decision. The Christian defenders, Greeks and Italians belonging to denominations that usually look down on one another, will enter Hagia Sophia in the city center for one final service together. Every candle beneath its great dome is lit, and the bells ring for a long time. An eyewitness later recalls the whole city heading toward the church, with even the lanes packed.

After the service, the emperor returns to the walls. Before dawn, the assault begins. After dawn, the defenses fall. The emperor dies amid the fighting, at a precise spot still unknown. Constantinople has fallen. The Byzantine Empire, the eastern half of the Roman Empire that had survived another thousand years on its own, ends here.

Remember this night. It occupies the same place in history textbooks around the world: printed at the midpoint of a timeline, with ``medieval'' on the left and ``early modern'' on the right. Many textbooks explicitly treat 1453 as a marker of the end of Europe's Middle Ages.

Now step outside that timeline and see what the rest of the world is doing on the same night.

In Beijing, this is Jingtai year four of the Ming dynasty. Four years earlier, the Oirats captured the Ming emperor Yingzong at Tumu. His younger brother Zhu Qiyu took the throne under the reign title Jingtai. Court politics seethe over whether and how to bring back the captured former emperor. Nobody in Beijing thinks the world has just entered a new age.

In France that July, French cannon defeat the English at Castillon, ending the Hundred Years' War after more than a century. For the French, 1453 is a year to celebrate, the opposite of the despair in Constantinople.

In Central America, the Aztecs expand temples in their lake capital, Tenochtitlan, while their ruler Moctezuma I extends his reach toward the Gulf of Mexico. In the South American Andes, the Inca Empire is expanding under Pachacuti; Machu Picchu's stone walls will rise block by block over the coming decades. In West Africa, Mali's caravans still carry salt and gold dust across the Sahara. People in these places have no concept of ``1453 CE.'' Each reckons years in its own way, each lives in its own present.

A timeline stamps this night as a boundary between eras. Yet for people in most of the world, it is an ordinary night, not even news. Word would take years to travel from Constantinople to Beijing or Tenochtitlan, if it ever arrived at all.

\section{Who Holds the Knife That Cuts Time?}
Let us put the conflict on the table without rushing to an answer.

History is too long to tell at once, so it must be divided. Nobody objects to that. The question is \textbf{where to cut}.

The familiar scheme divides history into antiquity, the Middle Ages, the early modern period, and the modern era. This scheme did not fall from the sky. Europeans fashioned it to cut their own history. The Western Roman Empire ended in 476, so antiquity ended; the Renaissance and oceanic discoveries arrived, so the early modern period began. Renaissance humanists themselves invented the term Middle Ages, meaning the period in between: between the classical age they admired and their own rebirth, an interval they regarded as a long wait.

It is a useful knife, reasonably suited to Europe's portion. Trouble began when it was used to cut the whole world. Chinese history was fitted into early antiquity, medieval times, and a later age; Indian, Islamic, and American histories were pushed into the same drawers. Awkward results followed. China's ``medieval'' period could run from Wei and Jin to Ming and Qing, stuffing more than a thousand years into one drawer. American history was sliced in two at 1492, with everything earlier called prehistory. Yet the Aztecs and Incas had writing or knotted-cord records, historical records, and immense cities. Their history was hardly ``pre.''

Another voice says: then stop cutting. History flows continuously, and periodization is just scholarly laziness. Satisfying as that sounds, it solves nothing. Without periods, you cannot even write a textbook's contents page. Without counting years, events three centuries apart merge into one. Medical students must make cuts in anatomy, and historians must make cuts in their accounts. The question is not whether to cut, but whether you know where you are cutting, why there, and how substantial the connecting tissue beneath the cut is.

This brings us to the chapter's real question:

\textbf{How should historical time be measured and divided? Whose divisions have the authority to place everyone in ``the same age''?}

If a Chinese person and a Spaniard stand in the year 1500, are they in the same era? Should the history of an Inca noble who has never heard of Rome contain a period called late antiquity? If periodization is merely convenient, what is convenience's price? Could we lose our way through our own time while holding someone else's map?

Before reading on, think for yourself. If you divided the time from humanity's emergence to today into several periods, which years would you cut at? Your choices reveal where you stand.

\section{One Year, Several Names}
First, listen to the positions people take on naming time. This is not a debate with one right and one wrong side, but several accounts, each with a full rationale from its own position.

\textbf{Voice one, the emperor: Time begins with my first year.}

For ancient Chinese dynasties, changing the reign title was a major political act. A new ruler would inaugurate an era; an auspicious sign could prompt another. A victory or a ceremony honoring Heaven and Earth could do so too. Emperor Wu of Han alone used eleven reign titles. The title worked somewhat like a brand today, announcing, ``The present belongs under my authority.'' Thus a coin bearing Kangxi tongbao does two jobs: it dates itself and reminds you who is Son of Heaven. Time enters power's ledger.

\textbf{Voice two, the farmer: Time is crops, not numbers.}

In the same Qing dynasty, an old farmer in a village a thousand li from the capital might not know which Kangxi year it was. His time moved like this: at Awakening of Insects, turn the soil; around Grain Rain, plant melons and beans; White Dew is too early for wheat, Cold Dew too late, but Autumn Equinox is just right. The twenty-four solar terms governed his rhythm. A new imperial reign title mattered far less than an early frost. This was not ignorance, but another equally precise timekeeping system, following the sun and seasonal signs to serve sowing and harvest. What we call the traditional agricultural calendar was never simply an emperor's calendar. It was another law of time.

\textbf{Voice three, the believer: Time begins with that event.}

The Islamic world counts from the Prophet Muhammad's migration to Medina, the Hijra, in 622 CE. Its calendar follows the moon, with about 354 days a year, so its New Year slowly passes through every season of the Gregorian calendar. The Jewish calendar counts from the traditional creation of the world and has now passed 5,700 years. The Buddhist era used in Thailand and elsewhere begins with Shakyamuni's passing into nirvana. Christians count from Jesus's birth. Every dating system is a declaration: for us, the world's most important event happened in that year. A calendar era's starting point is the brightest lamp in a community's memory.

\textbf{Voice four, the defeated: Your new age is my world's end.}

In 1492, Columbus's fleet, funded by the Spanish Crown, reached the Caribbean. European histories print the date large: the Age of Discovery, the beginning of the early modern world. That same year, Spain took Granada, ending the peninsula's last Muslim state, and soon afterward ordered the expulsion of Jews. For Indigenous Americans, the century after 1492 brought catastrophic population loss through war, enslavement, and diseases to which they lacked immunity. Many studies put the decline above half, and higher in some places. One date is labeled discovery, conquest, expulsion, and catastrophe by different people. The year itself is the same river, but those on different banks give it utterly different names.

\textbf{Voice five: A word in defense of dynastic chronologies.}

By now, you may want to tear the table of Chinese dynasties out of your textbook. Not yet. Let us state its case fully. This is an exercise worth taking seriously: present a position you disagree with so well that its supporters would nod in recognition.

The case has at least three layers. First, without a framework there is no history. Dynasties and dates are hooks. Hang major events on them, and there is somewhere to add details. Without any chronological framework, someone might imagine Yue Fei and Qi Jiguang as neighbors. Second, dynastic rise and fall is itself a real historical structure. Taxation, military recruitment, and eligibility for civil-service examinations really did follow dynastic changes. A chronology is not pure invention; it corresponds to the actual rhythm of political succession. Third, a shared chronology is a public language. Students across the country learn Tang, Song, Yuan, Ming, Qing, giving them minimum common coordinates for discussion. You cannot redefine time before every conversation.

This is not a weak case. The problem is never the chronology itself, but treating it as all of history, imagining that memorizing the hooks gives you everything hanging from them. A skeleton is necessary. Mistaking it for the whole body is the problem.

\section{Thinking Bridge: Turn the Timeline Sideways}
Here is a tool you can take away, designed for the problem of memorizing chronology without seeing history. I call it \textbf{parallel time bands}.

The method is simple. A textbook chronology runs vertically: one dynasty with its events, then the next with its own. Turn it sideways. Years run horizontally; the vertical axis lists regions rather than dynasties: China, India, the Islamic world, Europe, Africa, the Americas, one row each. Choose a year and slice down the rows to see what is happening everywhere at that moment.

Try a slice around 1500. On the top row, during the Ming Hongzhi reign, silk weaving thrives in Jiangnan's market towns and silver begins flowing in on a large scale. On another, Leonardo da Vinci has just finished The Last Supper, while Spain and Portugal discuss dividing the seas traversed by Columbus. Farther down, the Ottoman Empire's military power is formidable, inspiring both respect and fear among Venetian merchants. On India's row, the Mughal Empire does not yet exist, but its founder Babur is already sharpening his sword in Kabul. In the Americas, the Inca and Aztec empires are at their height, knowing nothing of the other side of the ocean. In Africa, Timbuktu in the Songhai Empire is West Africa's scholarly center, where scholars' book collections put most contemporary European libraries to shame.

A vertical list tells you that Qing follows Ming. Parallel bands tell you that these people were contemporaries. The two kinds of knowledge feel entirely different. The list creates the illusion of a relay race: one civilization finishes its run and passes the baton. The bands give a truer picture: everyone in history is crowded onto the same train, seated in different carriages and unable to hear one another.

This tool also corrects an easy mistake. We often equate different dynasties with people of different ages, imagining Qianlong as ancient and Washington as modern. Consider the facts: Qianlong died in the first lunar month of 1799, Washington that December. They were genuine contemporaries; Washington even died a few months later. Or consider that Oxford began teaching more than two hundred years before the Aztecs founded Tenochtitlan. The apparent mismatch reveals that our mental timelines for different regions and dynasties were drawn separately and never aligned. Parallel time bands force them into alignment.

Using the tool takes three steps: choose a year or decade, select three to five regions, and look up the people and events in each at that time, recording them along the same line. Repeat a few times and your questions change. Instead of asking why a dynasty fell, you ask why several places faced similar troubles in the same century. That is the kind of question historians find truly fascinating.

\section{Two Sources About ``Year One''}
Now read two short primary sources. The Spring and Autumn Annals, China's earliest surviving annalistic history, records 242 years of major events in Lu in a little over ten thousand Chinese characters. Its opening for the year Duke Yin took office consists of just six characters:

\begin{quote}
First year, spring, the king's first month. ---Spring and Autumn Annals, Duke Yin's First Year
\end{quote}
Literally: Duke Yin's first year, spring, the king's first month. Not one of the six characters records an event. They merely erect a temporal frame: which year, which season, whose calendar. But do not underestimate the final three characters. The king's first month means the first month proclaimed by the Zhou king. In theory, all the feudal states should honor the Zhou Son of Heaven's calendar, acknowledging his power to decide which month begins the year. Opening a historical work this way stamps its most conspicuous position with a declaration: time comes from the king. Later dynasties' practice of altering the year's beginning and issuing a calendar upon founding, treating calendrical authority as a standard sign of legitimate rule, has roots in this tradition. Whoever defines New Year's Day is the common ruler of the realm: one of China's earliest formulations of time as power.

Now the second source. In 221 BCE, Ying Zheng, king of Qin, unified the six states. Finding the title king insufficient, he invented a new one: huangdi, emperor. The ``Basic Annals of the First Emperor of Qin'' in the Records of the Grand Historian records his plans for what would follow him:

\begin{quote}
I shall be the First Emperor. Later generations shall be numbered, the Second, the Third, through the ten-thousandth generation, passing on without end.
\end{quote}
In other words: I am the First Emperor. My descendants will follow by number, Second, Third, all the way to Ten Thousand, forever. This was Chinese history's most ambitious plan for time: not only start counting with me, but continue to infinity. The result? The dynasty ended with the Second Emperor. From 221 to 207 BCE, it lasted only fifteen years counted inclusively. The ending is a lesson in itself. You can name time and plan its numbering, but time obeys nobody. Yet although Qin counted only to the Second Emperor, the imperial institution it created lasted more than two thousand years. Rupture and continuity often mingle. Which is the main thread depends on the distance from which you look.

Read the two sources together. One shows that naming time is a matter of power; the other that naming it cannot control time itself. These are the two anchors to which the rest of the chapter returns.

\section{476: Collapse or a Change of Form?}
Now practice comparing explanations. Take one of the best-known cuts in textbook chronology: in 476 CE, the Germanic commander Odoacer deposed the last Western Roman emperor. The standard wording is, ``The Western Roman Empire fell, and Western Europe entered the Middle Ages.'' Hidden in that statement is a judgment: a \textbf{rupture} occurred. Is rupture the only explanation? We can offer at least three, each with reasons and limits.

\textbf{Explanation one: Empire collapsed, civilization declined. The rupture view.}

Its case: the political map really shattered. The unified western Mediterranean broke into small Germanic kingdoms. Roman cities shrank, populations fell, the money economy receded, and much of Western Europe returned to barter. Great roads and aqueducts deteriorated. Latin literacy narrowed to small church circles. Archaeological evidence supports this dark picture. Affordable pottery common under Rome almost vanishes from Western European layers in the following centuries, and ordinary material life becomes noticeably rougher. This explanation's strength is respect for ordinary experience. A late-fifth-century Gallic farmer would not care about scholarly terminology, only that order had disappeared and tax collectors had become sword-bearing new masters.

Its limit: it speaks mainly from Rome's and the Roman aristocracy's viewpoint. For people in the eastern empire, the empire had not fallen at all. An emperor still sat in Constantinople and would for another millennium. For many farmers at the edges, the disappearance of imperial taxation might not be wholly bad news. When we say civilization collapsed, whose civilization do we mean?

\textbf{Explanation two: Rome did not die; it slowly changed clothes. The continuity view.}

Its case: the twentieth-century Belgian historian Henri Pirenne advanced a famous argument. Germanic kingdoms did not wish to destroy Rome; they admired it. Their kings legislated in Latin, retained Roman administrators, and valued recognition from Constantinople. What truly severed Mediterranean connections was the rise of the Arab world after the seventh century, turning Rome's inland sea into a boundary between worlds. Recent research goes farther with a concept called late antiquity. The centuries from the third to the eighth were not simply decline and fall, but a creative transition in which Christianity rose and culture gradually reorganized. This explanation reconnects the living institutions of the Middle Ages, Byzantium, the Church, monastic scribes, to Roman roots. History changes from a landslide into a river altering its course.

Its limit: too much continuity can leave us unable to speak of real suffering. Cities really became smaller, roads really broke down, and several generations really lived more dangerously. The word transition is too clean; it wipes away blood and mud. Explaining gradual change does not mean denying real falls within it.

\textbf{Explanation three: ``Rome's fall'' was a commemorative date invented later. The construction view.}

Its case: the event of 476 was rather unremarkable. The deposed emperor was a teenager. Odoacer did not kill him, but sent him to live in the countryside and continued writing to the emperor in Constantinople, presenting himself as merely administering Italy. Few contemporaries thought the sky had fallen, and substantial accounts are scarce. Later historians needing a dividing line selected and emphasized the symbol of a fall in 476. This reminds us that historical cuts often do not grow from events themselves. Historians insert them.

Its limit: pushed too far, this becomes nihilism. If every boundary is invented, history turns into an inseparable porridge. Yet from 550, looking back, Italy's political landscape really was utterly different from that of 400. The cut is chosen by people, but what is being cut is real.

Each explanation holds part of the truth. The exercise is not about selecting a correct answer, but developing a habit: \textbf{every periodization conceals an explanation, and every explanation brings a viewpoint and blind spots}. Next time you read that a certain year marks the beginning of an age, first ask: Who marked it, and from where? What continuities does this line erase?

Consider also an easily misjudged counterexample. A familiar story says that Constantinople fell in 1453, Greek scholars fled to Italy carrying ancient books, and the Renaissance began. The story is suspiciously smooth. Renaissance pioneers such as Petrarch were active in the fourteenth century, nearly a hundred years before 1453, and Greek scholars had already begun moving west before the city's fall. The fall did help carry texts westward. But hanging a cultural movement lasting two or three centuries on one city's capture is a classic shortcut of thinking that big events determine everything. It underestimates the slow currents below the surface, which take center stage next.

\section{Further Challenge: Time at Three Speeds}
Now for this chapter's weightiest theory. It comes from Fernand Braudel, a twentieth-century French historian of the Annales school. His major work studies the sixteenth-century Mediterranean. Yet instead of opening with rulers and generals, it devotes enormous space to mountains, seas, islands, trade routes, wheat, and olives. Many readers must wonder whether they bought history or geography.

Braudel did this deliberately. He believed historians had been viewing history at the wrong speed, and proposed that historical time flows simultaneously at at least three speeds.

\textbf{First, event time: fast, loud, like the news.} A battle, a coup, a treaty, an emperor's death. This is traditional history's staple, events lasting days or months. Braudel acknowledged its excitement but considered it superficial. In his famous metaphor, events are waves on history's surface.

\textbf{Second, conjunctural time: medium speed, measured in decades.} Population rises and falls, price cycles, shifting trade routes, an institution's growth and decay. Large silver inflows around the Ming--Qing transition changed China's economic fabric; after the Black Death, European labor became more expensive and slowly loosened the manorial system. One generation may not notice changes at this level. Two or three looking back can see them.

\textbf{Third, the longue durée: almost motionless, measured in centuries or millennia.} Geography, climate cycles, patterns of cultivation, deep technologies and ideas. A mountain determines which villages trade and which stay isolated. Monsoons set the rhythm of Indian Ocean commerce. The boundary between wheat and rice shapes differences between northern and southern life. These change very little, but lay out the tracks on which rapid change runs.

Why does this theory matter? It gives every question in this chapter a different foundation.

First consider what it explains. Why does the ancient, medieval, early modern scheme fit the world so awkwardly? Because it cuts according to event time: battles, dynasties, treaties. Yet the lives of an Indian farmer or a Saharan caravan guide are shaped by the second and third levels: monsoons, salt prices, crops, beliefs. Slow time does not obey dynasties. Once the layers are separated, the problem is clear. The world is not failing to keep up with European periods; the knife cuts only the surface.

This also explains why arguments about continuity and rupture never end. Rupture's advocates watch the first level, full of breaks. Continuity's advocates watch the second, full of gradual change. They are looking from different floors while imagining they stand at the same window.

But remain alert. Every grand theory leaves corners unlit, and Braudel's is no exception. The longue durée's greatest risk is writing people out of history. If monsoons, mountains, and population curves hold the steering wheel, what do particular choices count for? The fear and decisions of the soldier keeping watch on Constantinople's wall occupy not even a pixel in the longue durée. Yet history is made from countless such people's individual days. Recent historical work has visibly swung back toward individuals: a letter, a trial, a woman's lifetime account book, illuminating large structures from the smallest scale. This is called microhistory. Surface waves and deep ocean probably cannot replace each other.

A tool's limits are knowledge too. Braudel gives you glasses that switch among three focal lengths. Which to use remains your decision.

\section{Three Time Bands on Your Phone}
This old question is not old-fashioned at all.

First, news. The media's natural speed is Braudel's first level: trends measured in hours, headlines in days, shock and reversal as event time's standard sound effects. Learn about the world only through trends, and you watch the sea with a telescope that magnifies waves. You see them clearly but have no idea where the sea is flowing. Recognizing a news item's time level is a basic information-age defense. A sudden conflict may rest on thirty years of industrial change and three hundred years of geographical arrangements. Someone asking only what happened yesterday and someone saying only that things have been this way for three centuries will argue furiously because they are not discussing the same layer of time.

Next, consider how periodization governs your life. It is already divided into kindergarten, primary school, middle school, high school, university. Each has its uniforms, rules, and sense of an age, with solemn opening and graduation ceremonies between them, like small dynasties. A change of homeroom teacher resembles a new reign. This seems so natural that few ask why adolescents worldwide should be placed in similar compartments at similar ages. Educational periodization is a product of the last two centuries of industrial society. Its uniformity serves large-scale administration, not every individual's growth. Some people's Renaissance arrives early; others begin their Age of Discovery at forty. Periodization is a tool, not destiny.

Online, new eras are invented daily: year one of the internet, of short video, of the metaverse. Each is a small First Emperor-style operation, declaring that time begins here. Businesses invent first years to sell things; platforms invent birth-decade labels to create talking points. Every young generation is told it differs from all generations before, a statement every generation has heard once.

Most thought-provoking is the modern version of parallel time bands. Technology now makes worldwide presence at one moment possible: billions watch a match together; a message circles the globe in seconds. Humanity appears, for the first time, to share one timeline. But examine your screen more closely. On the evening you scroll through funny videos, someone your age elsewhere may be in a trench or queuing for water in a refugee camp. Technology aligns time, not circumstances. The question from that night in 1453 still glows in your pocket in another form. What happens elsewhere in the world at this same moment determines whether you can truly understand the present.

\section{Over to You: Where Would You Cut?}
Return to the opening coin.

Kangxi tongbao cast time into an emperor's numbering. Textbooks divide it into ancient, medieval, and early modern. Braudel separates three speeds. Your life is being divided into primary, middle, and high school. We have learned that all periodizations are human-drawn maps and every map leaves things out. We have also learned that without maps, movement is difficult.

The question now returns to your hands.

If you drew a timeline of your life from your earliest memories to today, where would you cut it into periods?

At events: a move, an exam, someone leaving? Or at slower changes: when you began looking at your parents differently, the year you suddenly understood a kind of conversation you had never understood before? In your chronology, which day divides before and after your era's starting point? Is the timeline others, your parents or teachers, draw for you the same as your own?

Think on a larger scale. A century from now, students will open history books and find our years assigned to an era with a name. Their cuts will probably miss something, erase some continuities, exaggerate some waves. Which date would you want them to remember, and which slow change would you fear they might forget?

There is no standard answer. But from your first cut, you are no longer merely a reader of time. You become one of its authors. Perhaps the whole secret of historical craft lies in that cut: it slices the past, but reveals where you stand.
